\documentclass{article}
\usepackage{amsmath}
\usepackage{hyperref}

\title{Comparison: EC50\_L Threshold Calculation Methods}
\author{}
\date{}

\begin{document}
\maketitle

\section*{Question}
Do \texttt{EC50\_lzd\_threshold\_sweep.py} and \texttt{MC\_ec50\_lzd.py} compute the
linezolid resistance-escape threshold ($EC_{50,L}$) the same way?

\section*{Conclusion}
Both scripts define and locate the threshold identically in principle; they differ
only in scope and robustness.

\subsection*{Shared methodology}
\begin{itemize}
  \item Identical \texttt{BASE\_PARAMS}, initial conditions $Y_0$, and ODE solver
        settings:
        \[
          \texttt{odeint(rtol=1e-7,\ atol=1e-9,\ mxstep=5000)}
        \]
  \item Threshold defined as the value of $EC_{50,L}$ at which the
        \emph{final} resistant reservoir count, $R_{res}(t_{\text{end}})$,
        crosses the limit of detection:
        \[
          f(EC_{50,L}) = R_{res}(t_{\text{end}};\,EC_{50,L}) - \mathrm{LOD}, \qquad \mathrm{LOD} = 10\ \text{CFU/mL}
        \]
  \item Root found via bisection using \texttt{scipy.optimize.brentq} over
        $EC_{50,L} \in [0.3,\,2.0]$ mg/L.
  \item Both scripts agree numerically: threshold $\approx 0.932$ mg/L at
        $\rho_S = 0.61$.
\end{itemize}

\subsection*{Differences}

\begin{enumerate}
  \item \textbf{Compartments bisected.}
  \begin{itemize}
    \item \texttt{EC50\_lzd\_threshold\_sweep.py} bisects \emph{both}
          $R_b$ (blood, index 1) and $R_{res}$ (reservoir, index 3)
          separately via \texttt{find\_threshold(compartment\_idx)}.
    \item \texttt{MC\_ec50\_lzd.py} bisects only $R_{res}$, producing a
          single \texttt{ESCAPE\_THRESH} used to classify Monte Carlo
          samples.
  \end{itemize}

  \item \textbf{Sweep resolution (diagnostic only).}
  \begin{itemize}
    \item \texttt{threshold\_sweep.py}: 171 points, step $0.01$ mg/L.
    \item \texttt{MC\_ec50\_lzd.py}: 121 points, step $\approx 0.014$ mg/L.
  \end{itemize}
  These grids feed only the visualization; \texttt{brentq} itself is
  continuous and unaffected by grid spacing.

  \item \textbf{No-crossing fallback behavior.}
  \begin{itemize}
    \item \texttt{threshold\_sweep.py}: if $f(EC50_{\min}) > 0$ or
          $f(EC50_{\max}) < 0$ (no sign change), returns \texttt{None}.
    \item \texttt{MC\_ec50\_lzd.py}: in the same situation, falls back to a
          coarse grid-based estimate,
          \[
            \texttt{ESCAPE\_THRESH} = \texttt{sweep\_ec50[argmax(sweep\_R\_res > LOD)]}
          \]
          so the Monte Carlo run does not fail outright.
  \end{itemize}
\end{enumerate}

\subsection*{Interpretation}
\texttt{EC50\_lzd\_threshold\_sweep.py} functions as the more rigorous
diagnostic tool (dual-compartment thresholds, finer sweep, strict
no-crossing handling), while \texttt{MC\_ec50\_lzd.py} only needs a single
$R_{res}$ threshold to partition its 1000-sample Monte Carlo population
into ``suppressed'' vs. ``escape'' outcomes, favoring robustness
(graceful fallback) over strictness.

\end{document}
